\documentclass[french,a4paper]{article}

\usepackage[utf8]{inputenc}
\usepackage[T1]{fontenc}
\usepackage{lmodern}
\usepackage{geometry}
\usepackage{graphicx}
\usepackage{babel}
\usepackage{amsmath}
\usepackage{amsthm}
\usepackage{amssymb}
\usepackage{hyperref}
\usepackage{stmaryrd}
\usepackage{enumitem}
\usepackage{tcolorbox}
\usepackage{dsfont}
\usepackage{multirow}
\usepackage{etoc}
\usepackage{titlesec}
\usepackage{esint}
\usepackage{cancel}
\usepackage{mathdots}

\setlist[itemize]{label=$\bullet$}


\renewcommand{\P}{\mathbb{P}}
\newcommand{\N}{\mathbb{N}}
\newcommand{\Z}{\mathbb{Z}}
\newcommand{\Q}{\mathbb{Q}}
\newcommand{\R}{\mathbb{R}}
\newcommand{\C}{\mathbb{C}}
\newcommand{\K}{\mathbb{K}}
\renewcommand{\L}{\mathbb{L}}
\newcommand{\U}{\mathbb{U}}
\newcommand{\st}{^*}
\newcommand{\tq}{\middle|}
\newcommand{\inv}{^{-1}}
\newcommand{\E}{\mathbb{E}}
\newcommand{\F}{\mathbb{F}}
\newcommand{\V}{\mathbb{V}}
\renewcommand{\ker}{\text{Ker}}
\newcommand{\im}{\text{Im}}
\newcommand{\card}{\text{Card}}
\newcommand{\Tr}{\text{Tr}}

\theoremstyle{plain}
\newtheorem{thm}{Théorème}
\newtheorem*{ind}{Indication}

\theoremstyle{definition}
\newtheorem{dfn}{Definition}

\theoremstyle{remark}
\newtheorem*{rmq}{Remarque}
\newtheorem*{app}{Application}

\newtcolorbox[auto counter]{ex}[2][]{colback=white, colframe=green!50!white, coltitle=black, fonttitle=\bfseries, title=Exercice~\thetcbcounter~: #2}

\title{TD Algèbre générale et arithmétique}

\begin{document}

\maketitle

\section{Groupes}

\begin{ex}{Centre d'un groupe}
	Soit $G$ un groupe. On note $Z(G)$ le centre de $G$ défini par $$Z(G)=\{x\in G\ \mid \ \forall y\in G,\ xy=yx\}$$ et, pour $x\in G$, on note : $$C(x)=\{gxg\inv\ \mid\ g\in G\}$$
	\begin{enumerate}
		\item Montrer que $Z(G)$ est un sous-groupes de $G$.
		\item Caractériser, à l'aide de $C(x)$, l'appartenance de $x$ à $Z(G)$.
	\end{enumerate}
\end{ex}

\begin{ex}{$\circ$}
	Dans un groupe, montrer que si $(ab)^n=e$, alors $(ba)^n=e$.
\end{ex}

\begin{ex}{Groupes de $\mathcal{M}_n(\K)$}
	Soit $n\in\N\st$. On appelle groupe de $\mathcal{M}_n(\K)$ toute partie de $\mathcal{M}_n(\K)$ stable par multiplication et qui, munie de la multiplication induite, est un groupe.
	\begin{enumerate}
		\item Le groupe $G$ est-il nécessairement un sous-groupe de $\mathcal{M}_n(\K)$ ?
		\item Montrer l'existence d'une unique matrice de projecteur dans $G$. On note $r$ son rang.
		\item $\star$ Montrer que $G$ est isomorphe à un sous-groupe de $\mathcal{GL}_r(\K)$.		
	\end{enumerate}
\end{ex}

\begin{ex}{Inversibles de $\Z/8\Z$}
	Déterminer le groupe des inversibles de $\Z/8\Z$. Ce groupe est-il cyclique ?
\end{ex}

\begin{ex}{CNS sur la cyclicité du produit}
	Donner une condition nécessaire et suffisante sur deux groupes cycliques $G$ et $H$ pour que le groupe produit $G\times H$ soit cyclique.
\end{ex}

\begin{ex}{Groupe des caractères}
	Soit $G$ un groupe cyclique d'ordre $n$. Déterminer $\hat{G}$ l'ensemble des morphismes de groupes de $G$ dans $(\C\st,\times)$. Montrer que $(\hat{G},\times)$ est cyclique.
\end{ex}

\begin{ex}{Racines $n$-èmes de l'unité}
	Pour tout $n\in\N\st$, on note $\U_n$ le groupe des racines $n$-èmes de l'unité. Déterminer $\U_m\cap\U_n$ et $\U_m\U_n$ pour $(m,n)\in(\N\st)^2$.
\end{ex}

\begin{ex}{Automorphismes de $\Z/p\Z$}
	\begin{enumerate}
		\item Déterminer les automorphismes du groupe $\Z/p\Z$ pour $p$ premier.
		\item En admettant que $(\Z/q\Z)\st$ est premier pour $q$ premier, déterminer l'ensemble des morphismes de groupes de $\Z/p\Z$ dans $\text{Aut}(\Z/q\Z)$ pour $p$ et $q$ premiers.
	\end{enumerate}
\end{ex}

\begin{ex}{$\Delta$ Sous-groupes additifs de $\R$ : HP à connaître}
	Les sous-groupes additifs de $\R$ sont soit de la forme $a\Z$ avec $a\in\R_+$, soit denses dans $\R$.
\end{ex}

\begin{proof}
	Soit $H$ un sous-groupe additif de $\R$. Si $H=\left\{0\right\}$, alors il est bien de la forme $a\Z$ avec $a=0$. Sinon, il contient un réel non nul, et quitte à passer à l'opposé, il contient un réel strictement positif. Donc $H\cap\R_+\st$ est non vide, et il admet une borne inférieure $a$.
	\begin{itemize}
		\item Supposons que $a\in H$, et montrons alors que $H=a\Z$. \\ D'une part, on a $<a>=a\Z\subset H$ par définition du sous-groupe engendré. Réciproquement, soit $x\in H$. Comme $a\in\R_+\st$, on peut considérer $q=\displaystyle\left\lfloor\frac{x}{a}\right\rfloor\in\Z$ et $r=x-aq$. Alors, d'après la première inclusion, $r\in H$, et comme $r\in[0,a[$, on a $r=0$ par minimalité de $a$. Donc $x=aq\in a\Z$.
		\item Supposons que $a\notin H\cap\R_+\st$., et montrons alors que $H$ est dense dans $\R$. \\ Donnons-nous un intervalle $]x,y[$ de largeur $\varepsilon=y-x>0$.
		\begin{itemize}
			\item Montrons qu'on peut trouver un $h\in]0,\varepsilon[$ dans $H$. Prenons dans $H\cap\R_+\st$ un $h_1$ tel que $h_1\in[a,a+\varepsilon[$. Comme $a\notin H\cap\R_+\st$, on a même $h_1\in]a,a+\varepsilon[$. Par le même raisonnement, prenons dans $H$ un $h_2\in]a,h_1[$. il suffit alors de poser $h=h_1-h_2$.
			\item Ensuite, posons $q=\displaystyle\left\lfloor\frac{x}{h}\right\rfloor+1\in\Z$. Alors $qh$ appartient à $H$ et il vérifie $x<qh<y$.
		\end{itemize}
		\item Enfin, aucun $H=a\Z$ ne peut être dense dans $\R$ (il s'agit bien d'un "soit/soit"). En effet :
		\begin{itemize}
			\item Si $a=0$, alors l'intervalle $]0,1[$ ne contient aucun élément de $H$.
			\item Si $a>0$, alors l'intervalle $]0,a[$ ne contient aucun élément de $H$.
		\end{itemize}
	\end{itemize}
	Ainsi, la preuve est achevée.
\end{proof}

\begin{ex}{$\star$ Sous-groupes fermés de $\C\st$}
	Déterminer les sous-groupes de $(\C\st,\times)$ qui sont fermés dans $\C$. On pourra utiliser librement le résultat classique sur les sous-groupes de $\R$.
\end{ex}

\begin{ex}{Vers les équations de Pell-Fermat}
	Soit $G=\left\{x+y\sqrt{2}\ \mid \ (x,y)\in\N\st\times\Z\ : \ x^2-2y^2=1\right\}$
	\begin{enumerate}
		\item Montrer que $G$ est un sous-groupe de $(\R_+\st,\times)$.
		\item $\star$ Montrer que $G$ est monogène.
	\end{enumerate}
\end{ex}

\begin{ex}{Une propriété sur les cardinaux de parties d'un groupe}
	Soit $A$ et $B$ deux parties d'un groupe fini $G$ telles que $\card(A)+\card(B)>\card(G)$. Montrer que $G=AB$ avec $$AB=\left\{ab \ |\ (a,b)\in A\times B\right\}$$
\end{ex}

\begin{ind}
	Pour l'inclusion directe, se donner $x \in G$ puis on pourra montrer que $A\inv x \cap B \ne \emptyset$.
\end{ind}

\begin{ind}
	Relativement à l'indication qui précède, on pourra observer que $A\cap B \ne \emptyset$ puis remarquer que $\card(A\inv x) = \card(A)$.
\end{ind}

\begin{proof}
	L'inclusion réciproque est immédiate par stabilité de $G$ en tant que groupe. Passons à l'inclusion directe. Déjà, on peut remarquer que $A\cap B \ne \emptyset$ : en effet, on a $A\cup B \subset G$ donc $$\card(A\cup B)=\card(A)+\card(B)-\card(A\cap B) \leqslant\card(G)$$ Par hypothèse, il s'ensuit que $$\card(A\cap B) \geqslant \card(A)+\card(B)-\card(G)>0$$ donc on a bien $A\cap B \ne \emptyset$. Soit désormais $x \in G$. L'application $$
	\begin{array}[t]{lrcl}
		f : & A & \longrightarrow & A\inv x \\
		& a & \longmapsto & a\inv x
	\end{array}
	$$ est bijective. En effet, elle est injective par régularité de $x$ et surjective par définition. Ainsi, on a $\card(A\inv x)=\card(A)$. Ensuite, on a donc l'inégalité stricte $\card(A\inv x)+\card(B) > \card(G)$. Exactement de la même façon que dans notre remarque introductive, on montre que $A\inv x \cap B \ne \emptyset$. Ainsi, on peut fixer $b \in A\inv x \cap B$. On peut alors fixer $a \in A$ tel que $b=a\inv x$. On obtient donc $x=ab \in AB$, ce qui achève l'exercice.
\end{proof}

\begin{ex}{Parties sans somme d'un groupe}
	Une partie $A$ d'une groupe abélien $(G,+)$ est dite \textbf{sans somme} lorsque : $$\forall (x,y)\in A^2,\ x+y\notin A$$
	\begin{enumerate}
		\item Soit $p$ un nombre premier congru à $2$ modulo $3$. On écrit $p=3k+2$.
		\begin{itemize}
			\item Montrer que l'ensemble $C$ des classes modulo $p$ des éléments de $\llbracket k+1,2k+1\rrbracket$ est une partie sans somme de $\Z/p\Z$.
			\item $\star$ Soit $A$ une partie non vide de $\Z/p\Z\setminus\{0\}$. Montrer qu'il existe un élément non nul $t$ de $\Z/p\Z$ tel que $$\card((tC)\cap A)\geqslant\frac{\card(A)}{3}$$ Indication : on pourra effectuer un raisonnement probabiliste.
		\end{itemize}
		\item Montrer qu'il existe une infinité de nombres premiers congrus à $2$ modulo $3$.
		\item Déduire de ce qui précède que dans toute partie finie non vide $A$ de $\N\st$, il existe une partie $B$ de $A$ sans somme et telle que $\displaystyle\card(B)\geqslant\frac{\card(A)}{3}$.
	\end{enumerate}
\end{ex}

\begin{ex}{$\Delta\star$ Une mesure de non commutativité}
	Soit $G$ un groupe fini non abélien. Montrer que $$\card\left(\left\{(x,y)\in G^2\ : \ xy=yx\right\}\right)\leqslant\frac{5}{8}\card(G)^2$$ et que l'égalité a lieu pour le groupe diédral $D_4$. Indication : on pourra examiner la structure de $C(x)=\left\{y\in G\ : \ xy=yx\right\}$ pour un $x\in G$ bien choisi.
\end{ex}

\begin{ex}{$\Delta$ Groupes ayant un nombre fini de sous-groupes}
	Caractériser les groupes qui n'ont qu'un nombre fini de sous-groupes.
\end{ex}

\begin{proof}
	Ce sont les groupes finis. Pour un tel groupe, tout élément est d'ordre fini, sinon il contient un sous-groupe isomorphe à $\Z$, puis on passe par les sous-groupes monogènes.
\end{proof}

\begin{ex}{Groupes diédraux}
	Pour $n\geqslant 3$, on note $D_n$ le groupes des isométries vectorielles du plan euclidien qui conservent un polygone régulier à $n$ côtés et d'isobarycentre $O$.
	\begin{enumerate}
		\item Le groupe $D_n$ est-il monogène ?
		\item Montrer qu'il admet une partie génératrice à deux éléments.
		\item Déterminer les classes de conjugaison des réflexions dans le groupe $D_n$.
	\end{enumerate}
\end{ex}

\begin{ex}{$\star$}
	Soit $m$ et $n$ dans $\N\st$. Déterminer tous les morphismes de $(\mathcal{GL}_n(\R),\times)$ dans $(\Z/m\Z,+)$.
\end{ex}

\begin{ex}{Produit semi-direct et cardinaux}
	Soit $G$ un groupe fini ainsi que $H$ et $K$ des sous-groupes de $G$. On définit : $$HK=\{hk\ \mid \ (h,k)\in H\times K\}$$ Montrer que $\card(HK)\times\card(H\cap K)=\card(H)\times\card(K)$. Attention : $HK$ n'est pas un sous-groupe en général !
\end{ex}

\begin{proof}
	Considérer l'application : $$\begin{array}{lrcl}
		f:&H\times K&\longrightarrow&G \\
		&(h,k)&\longmapsto&hk
	\end{array}$$ Remarquer que son image est exactement $HK$ et observer que tout élément de $HK$ admet exactement $\card(H\cap K)$ antécédents. Conclure en utilisant le lemme des bergers.
\end{proof}

\section{Ordre des éléments d'un groupe}

\begin{ex}{$\circ$}
	Montrer que tout groupe fini d'ordre pair contient un élément $x\ne e$ tel que $x^2=e$. Et pour les groupes d'ordre impair ?
\end{ex}

\begin{ex}{Un groupe dont tous les éléments sont d'ordre $2$ est commutatif}
	Montrer qu'un groupe dont tous les éléments sont d'ordre $2$ est commutatif.
\end{ex}

\begin{proof}
	Supposons $$\forall x\in G,\ x^2=e$$ Alors : $\forall x\in G,\ x=x\inv$. En appliquant ceci à $xy$, on obtient : $$xy=(xy)\inv=y\inv x\inv=yx$$ donc $G$ est commutatif.
\end{proof}

\begin{ex}{Groupe dont tous les éléments sont d'ordre 2}
	Soit $G$ un groupe d'élément neutre $e$ tel que $\forall x\in G,\ x^2=e$.
	\begin{enumerate}
		\item $\Delta$ Montrer que $G$ est abélien.
		\item Soit $H$ un sous-groupe de $G$ différent de $G$. Si $x$ est un élément de $G$ n'appartenant pas à $H$, montrer que $H\cup xH$ est un sous-groupe de $G$ isomorphe à $H\times \Z/2\Z$.
		\item Si $G$ est fini, que peut-on en déduire sur son cardinal et sa structure ?
		\item Donner un exemple d'un tel groupe infini.
	\end{enumerate}
\end{ex}

\begin{ex}{Petit théorème de Fermat sur les groupes : cas commutatif}
	Soit $G$ un groupe commutatif fini de cardinal $|G|$ et d'élément neutre $e$. Soit $x\in G$. En calculant $\displaystyle \prod_{g \in G}(xg)$ de deux façons différentes, montrer que $x^{|G|}=e$.
\end{ex}

\begin{proof}
	Déjà, remarquons que c'est la commutativité de $G$ qui nous autorise à définir ce produit correctement. Notons-le $P$ dans la suite. 
	\begin{itemize}
		\item D'une part, par commutativité, on a $P=x^{|G|}\displaystyle \prod_{g \in G}g$.
		\item D'autre part, l'application 
		$$\begin{array}[t]{lrcl} 
			f_x : & G & \longrightarrow& G \\
			& g & \longmapsto &xg
		\end{array}$$ 
		est une bijection. En effet, elle est injective par régularité de $x$ dans $G$, et elle est surjective car pour atteindre un élément $g$ de $G$, il suffit de prendre comme antécédent $x\inv g$. On peut donc effectuer un changement de variable qui prouve que $P=\displaystyle \prod_{g \in G}g$.
		\item Enfin, en égalisant ces deux expressions de $P$ et par régularité de $\displaystyle \prod_{g \in G}g$, on obtient bien que $x^{|G|}=e$.
	\end{itemize}
\end{proof}

\begin{ex}{$\Delta$ Petit théorème de Fermat pour les groupes : cas général}
	Soit $G$ un groupe fini (non nécessairement commutatif) de cardinal $|G|$ et d'élément neutre $e$. Soit $x \in G$. On souhaite montrer que $x^{|G|}=e$.
	\begin{enumerate}
		\item Justifier l'existence de $p=\min\left\{n \in\N\st \ | \ x^n=e\right\}$.
		\item Montrer que $\left\langle x\right\rangle=\left\{x^0, \ ..., x^{p-1}\right\}$ et en déduire que $\card(\left\langle x\right\rangle)=p$.
		\item Conclure.
	\end{enumerate}
\end{ex}

\begin{proof}
	Pour la 1, utiliser le principe des tiroirs.Pour la 2, utiliser la caractérisation du sous-groupe engendré par un élément et une division euclidienne. Montrer que les $x^k$ sont deux à deux distincts pour $k$ variant entre 0 et $p-1$. Pour la 3, utiliser le théorème de Lagrange (à savoir redémontrer car HP) ou utiliser sa variante au programme
	\begin{enumerate}
		\item Considérons 
		$\begin{array}[t]{lrcl}
			f_x : & \llbracket1, \ |G|+1\rrbracket & \longrightarrow& G \\
			& t & \longmapsto &x^t
		\end{array} $. D'après le principe des tiroirs, $f_x$ n'est pas injective. On peut donc fixer $1 \leqslant a,b \leqslant |G|+1$ tels que $a\ne b$ et $f_x(a)=f_x(b)$. On a alors $x^a=x^b$. Sans perte de généralité, quitte à échanger $a$ et $b$, on peut supposer que $a<b$. On a alors $x^{b-a}=e$ et $b-a \in \N\st$. L'ensemble que l'on considère est alors une partie de $\N$ non vide, donc son minimum existe.
		\item L'inclusion réciproque est immédiate car $\left\langle x\right\rangle=\left\{x^k \ | \ k \in\Z\right\}$. Dans l'autre sens, soit $y\in \left\langle x\right\rangle$. Fixons $k\in\Z$ tel que $y=x^k$. Effectuons la division euclidienne de $k$ par $p$ : $k=ap +b$ avec $a\in\Z$ et $b\in\llbracket0, \ p-1\rrbracket$. On a alors $$y=x^{ap+b}=(x^p)^ax^b=e^ax^b=x^b$$ et obtient bien le résultat.
		\\ \\ Pour montrer que le cardinal de cet ensemble vaut $p$, il suffit de montrer que ses éléments sont deux à deux distincts.  Soit $0 \leqslant m \leqslant n \leqslant p-1$ tels que $x^m=x^n$. Raisonnons par l'absurde : supposons que $n-m \geqslant p$. Alors $n\geqslant p+m \leqslant p > p-1$ ce qui est absurde. Donc $0 \leqslant n-m <p$. Or, $x^{n-m}=e$ donc par minimalité de $p$, $n-m=0$ puis $n=m$. Les éléments sont bien deux à deux distincts, donc on a bien $\card(\left\langle x\right\rangle)=p$.
		\item $\left\langle x\right\rangle$ est un sous-groupe de $G$, donc, d'après le théorème de Lagrange, son cardinal divise $|G|$. On peut donc fixer $k\in\Z$ tel que $|G|=kp$. On a alors $$x^{|G|}=x^{kp}=(x^p)^k=e^k=e$$ Pour mémoire, si on considère un sous-groupe $H$ quelconque de $G$ fini, il faut considérer la relation d'équivalence $\forall(x,y)\in G^2, \ x\sim y \Longleftrightarrow xy\inv \in H$. Toutes les classes ont même cardinal car ce sont les $xH$. Puisque les classes d'équivalences forment une partition, on obtient le théorème de Lagrange.
	\end{enumerate}
\end{proof}

\begin{ex}{Groupes d'ordre 6}
	Soit $G$ un groupe d'ordre $6$.
	\begin{enumerate}
		\item Montrer que $G$ possède au moins un élément d'ordre $3$ et au moins un élément d'ordre $2$.
		\item Montrer que si $G$ est commutatif, alors il est isomorphe à $\Z/6\Z$.
		\item On suppose $G$ non commutatif. On se donne dans $G$ des éléments $x$ et $y$ d'ordre respectifs $3$ et $2$. Montrer que l'application $$k\in\{0,1,2\}\mapsto x^kyx^{-k}$$ est injective. En déduire que $G$ possède exactement trois éléments d'ordre $2$, puis que $G$ est isomorphe au groupe symétrique $\mathcal{S}_3$.
	\end{enumerate}
\end{ex}

\begin{ex}{$\star$}
	Soit $G$ un groupe cyclique d'ordre $n$ et $d\in\N\st$.
	\begin{enumerate}
		\item Montrer que $G$ admet un sous-groupe d'ordre $d$ si, et seulement si, $d\mid n$, et qu'il y a alors unicité d'un tel sous-groupe.
		\item En déduire la formule d'Euler : $$n=\sum_{d\mid n}\varphi(d)$$
	\end{enumerate}
\end{ex}

\begin{ex}{Exposant d'un groupe}
	Montrer que si deux éléments ont des ordres premiers entre eux, alors leur produit est d'ordre fini. En déduire que tout groupe fini admet un élément d'ordre le ppcm de tous les ordres de ses éléments.
\end{ex}

\begin{ex}{Cyclicité}
	Soit $\K$ un corps. Montrer que tout sous-groupe fini $G$ de $\K\st$ est cyclique.
\end{ex}

\begin{rmq}
	En particulier, $(\Z/p\Z)\st$ est cyclique lorsque $p$ est premier.
\end{rmq}

\begin{ex}{$\Delta\star$ Lemme de Cauchy}
	Soit $G$ un groupe fini et $p$ un diviseur premier de $\card(G)$. On introduit : $$A=\left\{(g_1,\ldots,g_p)\in G^p\ : \ g_1\cdots g_p=1\right\}$$
	\begin{enumerate}
		\item Montrer que $(g_1,\ldots,g_p)\mapsto(g_1,\ldots,g_p,g_1)$ définit une permutation $\sigma$ de $A$.
		\item En considérant la décomposition de $\sigma$ en produit de cycles à support disjoint, montrer que $G$ possède un élément d'ordre $p$.
	\end{enumerate}
\end{ex}

\begin{ex}{Cyclicité des groupes d'ordre $pq$}
	Soit $p$ et $q$ des nombres premiers distincts. Montrer que tout groupe fini d'ordre $pq$ est cyclique. Indication : utiliser le lemme de Cauchy.
\end{ex}

\begin{proof}
	Avec le lemme de Cauchy, on dispose de $x\in G$ d'ordre $p$ et $y\in G$ d'ordre $q$. Comme $p$ et $q$ sont premiers entre eux car ce sont des nombres premiers distincts, on en déduit de façon assez classique que $xy$ est d'ordre $pq$. Comme $G$ possède un élément d'ordre son cardinal, $G$ est cyclique.
\end{proof}

\section{Groupe symétrique}

\begin{ex}{Conjugaison dans le groupe symétrique}
	Soit $E$ un ensemble. On note $\mathcal{S}(E)$ le groupe des permutations de $E$. On dit que deux éléments $\tau$ et $\tau'$ de $\mathcal{S}(E)$ sont conjugués s'il existe $\sigma\in\mathcal{S}(E)$ tel que $\tau'=\sigma\circ\tau\circ\sigma\inv$.
	\begin{enumerate}
		\item Montrer que deux cycles sont de même longueur si, et seulement s'ils sont conjugués.
		\item Application : en considérant les transpositions, déterminer tous les morphismes de groupe de $\mathcal{S}(\llbracket1,n\rrbracket)$ dans $\mathcal{C}\st$.
	\end{enumerate}
\end{ex}

\begin{ex}{Racine carré d'un $n$-cycle}
	Le cycle $(1\ 2\ \ldots\ n)$ admet-il une racine carrée dans $\mathcal{S}_n$ ?
\end{ex}

\begin{ex}{Un isomorphisme}
	Établir un isomorphisme entre le groupe des automorphismes de $\Z/2\Z\times\Z/2\Z$ et $\mathcal{S}_3$.
\end{ex}

\begin{ex}{Sous-groupes transitifs de $\mathcal{S}_n$}
	Un sous groupe $H$ de $\mathcal{S}_n$ est dit transitif lorsque $$\forall (x,y)\in\llbracket1,n\rrbracket^2,\ \exists\sigma\in H,\ \sigma(x)=y$$
	\begin{enumerate}
		\item Soit $\sigma\in\mathcal{S}_n$. A quelle condition le groupe $\langle\sigma\rangle$ engendré par $\sigma$ est-il transitif ?
		\item On considère une variable aléatoire $X_n$ de loi uniforme sur $\mathcal{S}_n$. Calculer la probabilité $p_n$ que le sous-groupe $\langle X_n\rangle$ engendré par $X_n$ soit transitif.
		\item On considère deux variables aléatoires indépendantes $X_n$ et $Y_n$ de même loi uniforme sur $\mathcal{S}_n$. On note $q_n$ la probabilité que le sous-groupe $\langle X_n,Y_n\rangle$ engendré par $X_n$ et $Y_n$ soit transitif; Montrer que $$q_n\xrightarrow[n\to+\infty]{}1$$
	\end{enumerate}
\end{ex}

\begin{ex}{$\star\star$ Nombre minimal de transpositions pour engendrer $\mathcal{S}_n$}
	Déterminer le nombre minimal de transpositions nécessaires pour engendre le groupe symétrique $\mathcal{S}_n$.
\end{ex}

\begin{proof}
	La réponse est $n+1$. On dispose de pleins de telles familles qui engendrent. Pour montrer qu'il en faut au moins $n+1$, utiliser des graphes.
\end{proof}

\begin{ex}{Matrices de permutations}
	Soit $n\in\N\st$. Pour $\sigma\in\mathcal{S}_n$, on note $P_\sigma=(\delta_{i,\sigma(j)})_{1\leqslant i,j\leqslant n}\in\mathcal{M}_n(\C)$.
	\begin{enumerate}
		\item Montrer que $\sigma\mapsto P_\sigma$ est un morphisme injectif de $\mathcal{S}_n$ dans $\mathcal{GL}_n(\C)$.
		\item $\star\star$ Soit $\sigma$ et $\tau$ dans $\mathcal{S}_n$. Montrer que $\sigma$ et $\tau$ sont conjuguées dans $\mathcal{S}_n$ si, et seulement si, les matrices $P_\sigma$ et $P_\tau$ sont semblables. Indication : on pourra étudier le polynôme caractéristique de $P_\sigma$.
	\end{enumerate}
\end{ex}

\begin{ex}{Nombre de permutations avec $k$ cycles distincts} On note $u_{n,k}$ le nombre de permutations de $\mathcal{S}_n$ avec $k$ cycles distincts (un point fixe compte pour un cycle). On a la formule de récurrence : $$u_{n+1,k+1}=u_{n,k}+n\times u_{n,k+1}$$
\end{ex}

\begin{proof}
	Le premier terme correspond au cas où $n+1$ est seul dans son cycle, et le deuxième aux cas où $n+1$ n'est pas seul dans son cycle.
\end{proof}

\section{Anneaux}

\begin{ex}{Caractéristique d'un anneau}
	Soit $A$ un anneau.
	\begin{enumerate}
		\item Montrer que $k\in\Z\mapsto k.1_A$ est l'unique morphisme d'anneaux de $\Z$ dans $A$.
		\item Montrer que le noyau de ce morphisme est de la forme $N\Z$ pour un unique $N\in\N$. $N$ est alors appelé \textbf{caractéristique de l'anneau $A$}.
		\item Que vaut $N$ lorsque :
		\begin{itemize}
			\item $A=\Z/n\Z$ pour $n\geqslant 2$ ;
			\item $A=\Z/n\Z\times\Z/m\Z$ pour $n\geqslant 2$ et $m\geqslant2$ ;
			\item $A=\R$ ;
			\item $A=\Z/n\Z\times \Z$ pour $n\geqslant 2$.
		\end{itemize}
		\item Montrer que si $A$ est un corps et $N\ne 0$, alors $N$ est un nombre premier et $$\forall (k,x)\in\Z\times A,\ k.x=0\iff (N\mid k\ \text{ou}\ x=0)$$
	\end{enumerate}
\end{ex}

\begin{ex}{$\circ$}
	Un groupe multiplicatif inclus dans un anneau $A$ est-il nécessairement un sous-groupe du groupe des inversibles de $A$ ? Et si $A$ est un corps ? intègre ?
\end{ex}

\begin{ex}{$\star$ Anneau noethérien}
	Soit $A$ un anneau. Montrer l'équivalence entre les assertions suivantes.
	\begin{enumerate}
		\item Tout idéal de $A$ est engendré par un nombre fini d'éléments.
		\item Toute suite croissante d'idéaux de $A$ est stationnaire.
		\item Toute famille non vide d'idéaux de $A$ possède un élément maximal.
	\end{enumerate}
	Lorsqu'un anneau vérifie l'une de ces trois conditions équivalentes, on dit que l'anneau est \textbf{noethérien}.
\end{ex}

\begin{rmq}
	On peut montrer qu'un anneau noethérien intègre est toujours un anneau factoriel, c'est-à-dire un anneau dans lequel a existence et unicité (à association près et à l'ordre des facteurs près) d'une décomposition en éléments irréductibles. Par exemple, tout anneau principal est noethérien puis factoriel, ce qui prouve par exemple de façon immédiate l'existence et l'unicité de la DFP dans $\Z$ et de la DFI dans $\K[X]$.
\end{rmq}

\begin{ex}{Idéaux premiers}
	Soit $I$ un idéal d'un anneau $A$. On dit que $I$ est premier lorsque : $$\forall (x,y)\in A^2,\ xy\in I\implies (x\in I\ \text{ou}\ y\in I)$$ Montrer que l'anneau quotient $A/I$ est intègre si, et seulement si, l'idéal $I$ est premier.
\end{ex}

\begin{ex}{Idéaux maximaux}
	Soit $I$ un idéal d'un anneau $A$. On dit que $I$ est maximal lorsque $I$ est un idéal propre de $A$ (inclus strictement dans $A$) et maximal au sens de l'inclusion. Montrer que l'anneau quotient $A/I$ est un corps si, et seulement si, l'idéal $I$ est maximal.
\end{ex}

\begin{ex}{Un anneau intègre non principal et non factoriel}
	Montrer que $\Z[i\sqrt{5}]$ est un anneau intègre. Déterminer l'ensemble de ses unités. Montrer que $$I=\{x+i\sqrt{5}y\ \mid (x,y)\in Z^2\ \text{et}\ 2\mid (y-x)\}$$ est un idéal. En supposant que cet idéal soit principal, montrer que le coefficient vaut $\pm 2$. En déduire une contradiction, et conclure que $A$ n'est pas principal. En considérant $6$, montrer que $\Z[i\sqrt{5}]$ n'est pas principal.
\end{ex}

\begin{ex}{$\circ$ Nilradical}
	Montrer que l'ensemble des éléments nilpotents d'un anneau commutatif $A$ est un idéal de $A$.
\end{ex}

\begin{ex}{Radical d'un idéal}
	Montrer que $$\sqrt{I}=\left\{x\in A \ \mid \exists n\in\N,\ x^n\in I\right\}$$ est un idéal de $A$ contenant $I$. Montrer que $\displaystyle\sqrt{\sqrt{I}}=\sqrt{I}$.
\end{ex}

\begin{ex}{Idéal engendré par les produits}
	Soient $I$ et $J$ deux idéaux d'un anneau commutatif $A$.
	\begin{enumerate}
		\item Montrer que le plus petit idéal de $A$ contenant tous les produits $ab$ pour $(a,b)\in A\times B$ est $$IJ=\left\{x\in A\ \mid\ \exists n\in\N,\ \exists((a_1,\ldots,a_n),(b_1,\ldots,b_n))\in I^n\times J^n,\ x=\sum_{i=1}^{n}a_ib_i\right\}$$ Le comparer à $I\cap J$.
		\item $\star$ A-t-on toujours $IJ=\{ab\ \mid (a,b)\in I\times J\}$ ?
	\end{enumerate}
\end{ex}

\begin{ex}{Idéaux maximaux dans un espace de fonctions}
	Montrer que les idéaux maximaux de $\mathcal{C}^0([0,1],\R)$ sont de la forme : $$I_x=\left\{f\in\mathcal{C}^0([0,1],\R)\ \mid\ f(x)=0\right\}$$
\end{ex}

\begin{proof}
	Je crois (à vérifier) qu'on peut utiliser le théorème de Borel-Lebesgue.
\end{proof}

\begin{ex}{Une réciproque du théorème chinois}
	Soit $a$ et $b$ deux entiers naturels non nuls. Montrer que les anneaux $(\Z/a\Z)\times(\Z/b\Z)$ sont isomorphes si, et seulement si, $a$ et $b$ sont premiers entre eux.
\end{ex}

\begin{ex}{Détermination de morphismes}
	Déterminer les morphismes d'anneaux, pus les morphismes de groupes, de $\Z/7\Z$ dans $\Z/13\Z$. Même question avec $\Z/6\Z$ et $\Z/15\Z$.
\end{ex}

\begin{ex}{Nilpotence dans $\Z/n\Z$}
	Déterminer une condition nécessaire et suffisante sur l'entier naturel pour qu'il y ait des éléments nilpotents non nuls dans $\Z/n\Z$.
\end{ex}

\begin{ex}{Anneau des entiers de Gauss}
	On pose $\Z[i]=\{x+iy\ \mid\ (x,y)\in\Z^2\}$.
	\begin{enumerate}
		\item Montrer que $\Z[i]$ est un sous-anneau de $\C$ stable par conjugaison.
		\item Montrer qu'un élément $u$ de l'anneau $\Z[i]$ est inversible si, et seulement si, $\lvert u\rvert=1$.
		\item Expliciter les éléments inversibles de $\Z[i]$.
	\end{enumerate}
\end{ex}

\begin{ex}{Anneau des entiers d'Eisenstein}
	On pose $j=\exp\left(2i\pi/3\right)$ et $\Z[j]=\{x+jy\ \mid\ (x,y)\in\Z^2\}$.
	\begin{enumerate}
		\item Montrer que $\Z[j]$ est un sous-anneau de $\C$ stable par conjugaison.
		\item Montrer qu'un élément $u$ de l'anneau $\Z[j]$ est inversible si, et seulement si, $\lvert u\rvert=1$.
		\item Expliciter les éléments inversibles de $\Z[j]$.
	\end{enumerate}
\end{ex}

\begin{ex}{Anneau des nombres décimaux}
	Montrer que l'ensemble des nombres décimaux $$\mathbb{D}=\{x\in\Q\ \mid \ \exists n\in\N,\ 10^nx\in\Z\}$$ est un anneau principal.
\end{ex}

\begin{ex}{Anneaux $\Z_p$}
	Soit $p$ un nombre premier et $\Z_p$ l'ensemble des rationnels de la forme $a/b$ avec $(a,b)\in\Z^2$ et $b\wedge p=1$.
	\begin{enumerate}
		\item Vérifier que $\Z_p$ est un sous-anneau de $\Q$.
		\item Quels sont ses éléments inversibles ? Ses éléments irréductibles ?
		\item Déterminer les idéaux de $A$. En déduire que $\Z_p$ est principal.
	\end{enumerate}
\end{ex}

\begin{ex}{$\star$ Principalité des sous-anneaux de $\Q$}
	Montrer que tout sous-anneau $A$ de $\Q$ est principal.
\end{ex}

\section{Corps}

\begin{ex}{$\circ$}
	Quels sont les corps $\K$ tels que $\forall x\in\K\st,\ x\inv=x$ ?
\end{ex}

\begin{ex}{$\Delta$ Injectivité des morphismes de corps}
	Montrer que tout morphisme de corps est injectif.
\end{ex}

\begin{proof}
	Soit $\varphi$ un morphisme du corps $\K$ vers le corps $\L$. Raisonnons par l'absurde : supposons que $\varphi$ n'est pas injectif. On peut alors fixer $x \in \ker(\varphi) \setminus \{0_\K\}$. Or, $x$ est non nul, donc on a : $$1_\L=\varphi(1_\K)=\varphi(xx\inv)=\varphi(x)\varphi(x\inv)=0_\L$$ par propriétés des morphismes de corps ainsi que par absorbance. Or, nos corps sont toujours supposés non triviaux, donc ceci est \textbf{absurde}. Ainsi, $\varphi$ est bien injectif, ce qui conclut.
\end{proof}

\begin{ex}{$\Delta$ Structure de corps des anneaux intègres finis}
	Montrer que tout anneau intègre fini est un corps.
\end{ex}

\begin{proof}
	Donnons-nous un anneau intègre fini $A$ et $a$ un élément non nul de $A$. Considérons l'application
	$$
	\begin{array}[t]{lrcl}
		\varphi_a : & A & \longrightarrow & A \\
		& x & \longmapsto & ax
	\end{array}
	$$
	Déjà, cette application est bien définie par stabilité de l'anneau par la deuxième loi. Ensuite, elle injective par intégrité de $A$. Enfin, elle est alors surjective par égalité de cardinaux. Donc on peut fixer un antécédent de $1_A$ par $\varphi_a$, ce qui montre que $a$ est inversible. Ainsi, $A$ est bien un corps car seule l'inversibilité des éléments non nuls manquait.
\end{proof}

\begin{rmq}
	On peut montrer de même que toute algèbre intègre de dimension finie est un corps.
\end{rmq}

\begin{ex}{Idéaux d'un corps}
	Déterminer l'ensemble des idéaux d'un corps $\K$
\end{ex}

\begin{ind}
	Que se passe-t-il si un élément non nul de $\K$ appartient à un idéal ?
\end{ind}

\begin{proof}
	On raisonne classiquement par analyse-synthèse.
	\begin{itemize}
		\item \textbf{Analyse} : Soit $I$ un idéal de $\K$. L'idéal nul est évidemment solution, donc on peut désormais supposer que $I$ n'est pas réduit au groupe trivial. Fixons $x \in I\setminus\left\{0_\K \right\}$. Puisque $I$ est un idéal, $xx\inv=1_\K \in I$. Par suite, pour tout $y \in\K$, $y\times 1_\K=y \in I$, donc $I=\K$.
		\item \textbf{Synthèse} : Réciproquement, $\left\{0_\K \right\}$ et $\K$ sont bien des idéaux de $\K$, la vérification est immédiate.
	\end{itemize}
\end{proof}

\begin{ex}{Cardinal d'un corps fini}
	Démontrer que tout corps fini $\K$ est de cardinal $p^n$ avec $p$ un nombre premier et $n\in\N\st$.
\end{ex}

\begin{proof}
	Considérer $$\L=\{k\mathbf{1}_\K \ \mid \ k\in\Z\}$$ qui est un sous-corps de $\K$. Démontrer que sa caractéristique est un nombre premier. Enfin, voir $\K$ comme un $\L$-espace vectoriel et conclure par dimension.
\end{proof}

\begin{ex}{Éléments algébriques}
	Soit $\L$ un sur-corps de $\K$. on dit que $x\in\L$ est algébrique sur $\K$ lorsqu'il admet un polynôme annulateur non nul, c'est-à-dire un polynôme $P\in\K[X]$ tel que $P(x)=0$.
	\begin{enumerate}
		\item Montrer que si $x\in\L$ est algébrique sur $\K$, alors il existe un polynôme annulateur $\Pi$ non nul de degré minimal et que celui-ci est irréductible. Montrer alors que $\K[x]$ est un corps de dimension $\deg(\Pi)$ comme $\K$-espace vectoriel.
		\item Montrer que $x\in\L$ est algébrique sur $\K$ si, et seulement si, $\K[x]$ est de dimension finie sur $\K$.
		\item $\star$ Montrer que l'ensemble des éléments algébriques sur $\K$ est un sous-corps de $\l$.
		\item $\star\star$ Montrer que l'ensemble des nombres complexes algébriques sur le corps $\Q$ est un corps algébriquement clos, c'est-à-dire dans lequel tout polynôme est scindé. Est-il isomorphe à $\C$ ?
	\end{enumerate}
\end{ex}

\section{Arithmétique des entiers}
\begin{ex}{Formule d'Euler}
	Montrer la formule d'Euler : $$n=\sum_{d\mid n}\varphi(d)$$ où $\varphi$ est l'indicatrice d'Euler.
\end{ex}
\begin{ex}{Produits d'entiers successifs et carrés parfaits}
	Trouver les $n\in\N\st$ tels que $n(n+1)$ soit un carré parfait, puis les $n\in\N\st$ tels que $n(n+1)(n+2)$ soit un carré parfait.
\end{ex}
\begin{ex}{Il y a de gros trous dans la liste !}
	Montrer qu'il existe des listes d'entiers consécutifs arbitrairement grandes ne contenant aucun nombre premier.
\end{ex}
\begin{proof}
	Considérer les nombres $n!+2$, $n!+3$, $\ldots$, $n!+n$.
\end{proof}
\begin{ex}{Nombres premiers congrus à $-1$ modulo $4$}
	Démontrer qu'il existe une infinité de nombres premiers congrus à $-1$ modulo $4$.
\end{ex}
\begin{ex}{Carrés dans $\F_p$ et nombres premiers congrus à $1$ modulo $4$}
	Soit $p$ un nombre premier impair.
	\begin{enumerate}
		\item Soit $a\in\F_p$ non nul. Montrer que $a^{(p-1)/2}=\pm 1$. Et si $a$ est un carré dans $\F_p$ ?
		\item Montrer que $a$ est un carré dans $\F_p$ si, et seulement si, $a^{(p-1)/2}=1$. Indication : on pourra considérer le produit des éléments non nuls de $\F_p$ ainsi que les deux involutions $x\mapsto x\inv$ et $x\mapsto ax\inv$ de $\F_p\setminus\{0\}$. Sinon, on pourra aussi dénombrer les carrés de $\F_p$ à l'aide d'un morphisme.
		\item En déduire que $-1$ est un carré dans $\F_p$ si, et seulement si, $p$ est congru à $1$ modulo $4$.
		\item Montrer qu'il existe une infinité de nombre premiers congrus à $1$ modulo $4$.
	\end{enumerate}
\end{ex}
\begin{ex}{$\star$}
	Soit $p$ un nombre premier impair et $a\in\Z$ non multiple de $p$. Montrer que s'il existe $k\in\Z$ tel que $k^2\equiv a[p]$ alors, pour tout $n\in\N\st$, il existe $k\in\Z$ tel que $k^2\equiv a[p^n]$. Généraliser à des polynômes autres que $X^2-a$.
\end{ex}
\begin{ex}{Une divisibilité étrange}
	Déterminer les entiers $n\in\N\st$ tels que $7$ divise $n^n-3$.
\end{ex}
\begin{ex}{Critères de divisibilité}
	Déterminer un critère de divisibilité par $31$. Généraliser.
\end{ex}
\begin{ex}{Une factorisation de Sophie Germain}
	Déterminer les entiers naturels $n$ pour lesquels $n^4+4$ est premier.
\end{ex}
\begin{proof}
	Écrire $n^4+4=n^4+4n^2+4-(2n)^2$ puis factoriser et discuter.
\end{proof}
\begin{ex}{Une équation du second degré dans $\Z/143\Z$}
	Trouver les racines modulo $143$ de l'équation $x^2+x+11=0$.
\end{ex}
\begin{proof}
	Remarquer que $143=11\times 13$ et utiliser le théorème chinois.
\end{proof}
\begin{ex}{Une non divisibilité}
	Montrer que $121$ ne divise jamais $n^2+3n+5$.
\end{ex}
\begin{ex}{Étrange...}
	Déterminer la somme des chiffres de la somme des chiffres de la somme des chiffres de $4444^{4444}$.
\end{ex}
\begin{proof}
	Passer modulo $9$ pour obtenir la congruence modulo $9$ du nombre cherché. Ensuite, majorer extrêmement brutalement les sommes des chiffres, ce qui fait énormément diminuer la valeur possible, puis conclure en observant que seule un valeur est possible dans l'encadrement finale avec la congruence modulo $9$.
\end{proof}
\begin{ex}{$\Delta$ Théorème de Wilson}
	Soit $n\in\N$ avec $n\geqslant 2$. Montrer que $n$ est premier si, et seulement si, $(n-1)!\ \equiv\ -1\ [n]$.
\end{ex}
\begin{ex}{$\star$}
	Soit $p$ un nombre premier et $q\in\N\st$ premier avec $p$. Pour tout $n\in\N\st$, on note $t_n$ l'ordre de $\overline{q}$ dans le groupe des inversibles de $\Z/p^n\Z$. Montrer que la suite $(t_{n+1}/t_n)_{n\geqslant 1}$ est constante à partir d'un certain rang. Montrer que le résultat reste vrai sans supposer $p$ premier.
\end{ex}
\begin{ex}{$\Delta$$\star$ Formule de Legendre}
	Soit $n\in\N\st$ et $p$ un nombre premier. Montrer la formule de Legendre : $$v_p(n!)=\sum_{i=1}^{+\infty}\left\lfloor\frac{n}{p^i}\right\rfloor$$
\end{ex}
\begin{proof}
	Voici une preuve très élégante due à François Moulin. On munit $\Omega=\llbracket1,n\rrbracket$ de la probabilité uniforme et on s'intéresse à la variable aléatoire $v_p$ définie sur $\Omega$, à valeurs dans $\N$. On dispose alors de la formule : $$\E(v_p)=\sum_{k=1}^{n}\frac{1}{n}v_p(k)$$ par la formule de transfert. Mais, par propriété de la valuation $p$-adique, on a alors : $$\E(v_p)=\frac{1}{n}v_p(n!)$$ D'autre part, puisque la variable aléatoire est à valeurs dans $\N$, on dispose aussi de la formule : $$\E(v_p)=\sum_{i=1}^{+\infty}\P(v_p\geqslant i)$$ Or, puisque l'on a choisi la probabilité uniforme, il suffit alors de déterminer le nombre d'entiers de $\llbracket1,n\rrbracket$ dont la valuation $p$-adique est supérieure à $i$, c'est-à-dire le nombre d'entiers de $\llbracket1,n\rrbracket$ divisibles par $p^i$. Or, il y en a exactement $\displaystyle\left\lfloor\frac{n}{p^i}\right\rfloor$. Ainsi : $$\E(v_p)=\sum_{i=1}^{+\infty}\frac{1}{n}\left\lfloor\frac{n}{p^i}\right\rfloor$$ En égalisant les deux expressions différentes de $\E(v_p)$ et en multipliant par $n$, on obtient finalement : $$v_p(n!)=\sum_{i=1}^{+\infty}\left\lfloor\frac{n}{p^i}\right\rfloor$$
	Il existe principalement deux autres preuves : la première, fondée sur le même principe du nombre d'entiers de $\llbracket1,n\rrbracket$ divisibles par $p^i$ suivie d'un télescopage, et une autre par récurrence en utilisant des propriétés de la partie entière. Néanmoins, la preuve proposée ici, très originale, a le mérite d'être très brève est assez simple.
\end{proof}
\begin{ex}{Convolution de Dirichlet}
	On note $E$ l'espace des fonctions de $\N\st$ dans $\C$. On le munit de la loi $*$, dite convolution de Dirichlet, définie par : $$\forall (f,g)\in E^2,\ \forall n\in\st,\ (f*g)(n)=\sum_{d\mid n}f(d)g\left(\frac{n}{d}\right)$$
	\begin{enumerate}
		\item Montrer que $E$ est un monoïde commutatif de neutre $\varepsilon:n\mapsto\delta_{1,n}$.
		\item Montrer que la convolution de deux fonctions arithmétiques multiplicatives est arithmétique multiplicative. On rappelle qu'une fonction $f$ est arithmétique multiplicative si elle est définie sur $\N\st$, si $f(1)=1$ et si $m\wedge n=1\implies f(mn)=f(m)f(n)$.
	\end{enumerate}
\end{ex}
\begin{ex}{$\Delta$ Fonction de Möbius}
	On définit la fonction de Möbius $\mu$ sur $\N\st$ de la façon suivante : à $n\in\st$, elle associe
	\begin{itemize} 
		\item $(-1)^r$ si $n$ est le produit d'exactement $p$ nombres premiers deux à deux distincts ;
		\item $0$ sinon. 
	\end{itemize}
	\begin{enumerate}
		\item Montrer que la fonction de Möbius est arithmétique multiplicative.
		\item Montrer que : $$\forall n\in\N\st,\ \sum_{d\mid n}\mu(d)=\delta_{1,n}$$
		\item En déduire, en utilisant l'exercice sur la convolution de Dirichlet, la formule d'inversion de Möbius : $$\forall(f,g)\in E^2, \left(\forall n\in\N\st,\ g(n)=\sum_{d\mid n}f(d)\right)\iff\left(\forall n\in\N\st,\ f(n)=\sum_{d\mid n}\mu(d)g\left(\frac{n}{d}\right)\right)$$
	\end{enumerate}
\end{ex}
\begin{rmq}
	En utilisant la formule d'Euler, on peut en déduire que que $$\forall n\in\N\st,\ \varphi(n)=\sum_{d\mid n}\mu\left(\frac{n}{d}\right)d$$
\end{rmq}
\begin{ex}{Autour du nombre de diviseurs et généralisation}
	Pour tout $k\in\N$, on définit $\sigma_k$ sur $\N\st$ par : $$\sigma_k:n\mapsto\sum_{d\mid n}d^k$$ où les $d$ qui divisent $n$ sont pris positifs.
	\begin{enumerate}
		\item Montrer que $\sigma_0(n)$ est le nombre de diviseurs de $n$.
		\item Soit $k\in\N$. Montrer que $\sigma_k$ est une fonction arithmétique multiplicative.
		\item Soit $k\in\N$ et $p$ un nombre premier. Montrer que $$\forall\alpha\in\N,\ \sigma_k(p^\alpha)=\sum_{i=0}^{\alpha}p^{ik}$$
		\item Si $n=p_1^{\alpha_1}\times\cdots\times p_r^{\alpha_r}$ est la décomposition en facteurs premiers de $n$ (les $p_i$ sont deux à deux distincts), montrer que $$\sigma_0(n)=\prod_{i=1}^{r}(1+\alpha_i)$$
		\item Plus généralement, pour tout $k\in\N$, sous les mêmes hypothèses qu'à la question précédente, montrer que $$\sigma_k(n)=\sum_{i=1}^{r}\prod_{j=1}^{\alpha_i}p_i^{jk}$$
	\end{enumerate}
\end{ex}
\begin{proof}
	\begin{enumerate}
		\item Utiliser la définition et compter.
		\item Utiliser le lemme de Gauss.
		\item Déterminer les diviseurs de $p^\alpha$.
		\item Utiliser la DFP pour trouver le nombre de diviseurs.
		\item Utiliser les questions 2 et 3.
	\end{enumerate}
\end{proof}
\begin{ex}{$\star$ Méthode de l'hyperbole de Dirichlet}
	Si $d$ est la fonction qui à $k$ associe son nombre de diviseurs positifs, alors on a : $$\sum_{k=1}^{n}d(k)=n\ln(n)+(2\gamma-1)n+O(\sqrt{n})$$
\end{ex}
\begin{ex}{$\star$}
	Soit $P\in\Z[X]$ non constant. Montrer qu'il existe une infinité de nombres premiers $p$ pour lesquels $p$ divise $P(n)$ pour au moins un entier $n\in\Z$. Indication : s'intéresser à $P(n+qP(n))$.
\end{ex}
\begin{proof}
	Utiliser l'indication et utiliser le fait que $b-a$ divise $P(b)-P(a)$ pour $(a,b)\in\Z^2$.
\end{proof}
\begin{ex}{$n$ ne divise pas $2^n-1$}
	Soit $n$ un entier supérieur ou égal à $2$. Montrer que $n$ ne divise pas $2^n-1$. Indication : on pourra considérer le plus petit diviseur premier de $n$.
\end{ex}
\begin{ex}{$\Delta$ Nombres de Mersenne}
	\begin{enumerate}
		\item Montrer : $a\mid n\implies 2^a-1\mid 2^n-1$.
		\item Condition nécessaire pour que $2^n$ soit premier.
		\item Montrer que les diviseurs premiers de $2^p-1$ avec $p$ premier impair sont de la forme $2kp+1$.
		\item $2^{11}-1$ est-il premier ?
	\end{enumerate}
\end{ex}

\begin{ex}{Lemme de Wolstenholme}
	Soit $p$ un nombre premier. On écrit le rationnel $\displaystyle\sum_{i=1}^{p-1}\frac{1}{i}$ sous forme irréductible $\displaystyle\frac{a}{b}$. Montrer que $p^2$ divise $a$. On s'intéressera au polynôme $Q=\displaystyle\prod_{k=1}^{p-1}(X-k)$ considéré dans $\Z[X]$ ou $\F_p[X]$.
\end{ex}

\section{TD Châteaux}
On notera la loi de groupe multiplicativement lorsque l'exercice ne traite que de groupes. Pour $n\in\N\st$, on note $\varphi(n)$ le nombre d'entiers inférieurs à $n$ et premiers avec $n$. On désigne par $(\Z/n\Z)^\times$ le groupe multiplicatif des inversibles de l'anneau $\Z/n\Z$. On rappelle que selon la définition du programme, tous les corps sont supposés commutatifs \textit{a priori}.
\subsection{Sous-groupes distingués et groupes quotients}
\begin{ex}{Sous-groupes distingués et groupes quotients}
	Soit $g\in G$. On note $\sigma_g$ l'automorphisme intérieur $x\mapsto gxg\inv$ de $G$.
	\begin{enumerate}
		\item Vérifier que $\sigma_g$ est un automorphisme de $G$.
		\item Vérifier que l'application $g\mapsto\sigma_g$ est un morphisme de $G$ dans $\text{Aut}(G)$, dont le noyau est le centre de $G$.
		\item On définit dans $G$ la relation binaire $\mathcal{R}$ par $x\mathcal{R}y\iff \exists g\in G,\ y=\sigma_g(x)$. On vérifie sans difficulté qu'il s'agit d'une relation d'équivalence. La classe de $x$ est appelée orbite de $x$, notée $\mathcal{O}_x$. A quelle condition l'orbite de $x$ est-elle réduite au singleton $\{x\}$ ? Soit $G$ un groupe et $H$ un sous-groupe de $G$. On définit la relation binaire $\mathcal{R}_g$ par : $$x\mathcal{R}_gy\iff x\inv y\in H$$ On définit de même la relation binaire $\mathcal{R}_d$ par $x\mathcal{R}_dy\iff xy\inv\in H$. On vérifie facilement que $\mathcal{R}_g$ est une relation d'équivalence et que la classe d'équivalence d'un élément $x$ de $G$ est égale à $xH$. On appelle indice de $H$ dans $G$ le nombre de classes d'équivalence selon $\mathcal{R}_g$ (s'il est fini). On le note $[G:H]$. De la même façon, $\mathcal{R}_d$ est aussi une relation d'équivalence et la classe d'équivalence de $x$ selon $\mathcal{R}_d$ est égale à $Hx$. Si $H$ est un sous-groupe de $G$, on dit que $H$ est distingué dans $G$ lorsque $\forall x\in H,\ \forall g\in G,\ gxg\inv\in H$. Si $H$ est un sous-groupe distingué de $G$, les relations d'équivalence $\mathcal{R}_g$ et $\mathcal{R}_d$ coïncident, la classe d'équivalence étant $xH$. L'ensemble quotient est alors noté $G/H$. On montre alors aisément que la relation $\mathcal{R}_g$ est compatible avec la loi de $G$ et la loi quotient muni $G/H$ d'une structure de groupe, appelé groupe quotient de $G$ par $H$. Si $G$ est fini, le groupe quotient est d'ordre $\lvert G\rvert/\lvert H\rvert$.
		\item Déterminer les sous-groupes distingués du groupe symétrique $\mathcal{S}_4$.
	\end{enumerate}
\end{ex}
\subsection{Actions de groupes et formule des classes}
\begin{ex}{Actions de groupes}
	Soient $E$ un ensemble et $G$ un groupe. une action du groupe $G$ sur l'ensemble $E$ est une application $$\begin{array}{rcl}
		G\times E\longrightarrow&E \\
		(g,x)&\longmapsto&g\cdot x
	\end{array}$$ qui vérifie $\forall (g,g')\in G^2,\ \forall x\in E,\ (gg')\cdot x=g\cdot(g'\cdot x)$ et $e\cdot x=x$. Cela revient à se donner pour tout $g\in G$ une permutation $\sigma_g:E\to E$ définie par $\sigma_g(x)=g\cdot x$, avec la propriété suivante : $\forall (g,g')\in G^2,\ \sigma_g\circ\sigma_{g'}=\sigma_{gg'}$ et $\sigma_e=\text{Id}$. Si $x\in E$, on définit l'orbite de $x$ comme étant $\mathcal{O}_x=\{g\cdot x\ \mid \ g\in G\}$.
	\begin{enumerate}
		\item Vérifier que l'application $$\begin{array}{rcl}
			G^2&\longrightarrow&G\\
			(g,x)&\longmapsto gx
		\end{array}$$ est une action de groupe, qu'on appelle action par translation.
		\item Même question avec l'application $$\begin{array}{rcl}
			G^2&\longrightarrow&G\\
			(g,x)&\longmapsto gxg\inv
		\end{array}$$ qu'on appelle action par automorphisme intérieur.
		\item Vérifier que les orbites forment une partition de $E$.
		\item Si $x\in E$, on définit le stabilisateur de $x$ comme étant $\text{Stab}_x=\{g\in G\ \mid \ g\cdot x=x\}$. Vérifier que $\text{Stab}_x$ est un sous-groupe de $G$.
		\item En appliquant le lemme des bergers, montrer que si $G$ est fini, pour tout $x\in E$, $\card(G)=\card(\mathcal{O}_x)\card(\text{Stab}_x)$.
	\end{enumerate}
\end{ex}
\begin{ex}{Équation aux classes}
	On se donne une action du groupe $G$ sur l'ensemble $E$. On pose $E_G=\{x\in E\ \mid\ \forall g\in G,\ g\cdot x=x\}$. Remarquons qu'alors : $x\in E_G\iff \mathcal{O}_x=\{x\}$.
	\begin{enumerate}
		\item Si $G$ est fini, montrer que $$\card(E)=\card(E_G)+\sum_{x\in A'}\card(\mathcal{O_x})$$ où $A'$ est un ensemble de représentants des orbites non réduites à un singleton.
		\item Si $E$ et $G$ sont finis, en déduire l'équation aux classes : $$\card(E)=\card(E_G)+\sum_{x\in A'}\frac{\card(\mathcal{O}_x)}{\card(\text{Stab}_x)}$$
		\item Appliquer la formule précédente à l'action de $G$ sur lui-même par automorphisme intérieur.
		\item Pour $g\in G$, on pose $E_g=\{x\in E\ \mid\ g\cdot x=x\}$. On note $r$ le nombre d'orbites de l'action de $G$ sur $E$. Démontrer la formule de Burnside : $$r=\frac{1}{\card(G)}\sum_{g\in G}\card(E_g)$$
	\end{enumerate}
\end{ex}
\begin{ex}{Application aux $p$-groupes}
	Soit $p$ un nombre premier.
	\begin{enumerate}
		\item Démontrer que tout groupe d'ordre $p^2$ est abélien.
		\item Démontrer que tout groupe d'ordre $p^k$ pour $k\geqslant 3$ possède un centre non réduit à l'élément neutre.
	\end{enumerate}
\end{ex}
\subsection{Autour du théorème de Cauchy}
Dans les exercices qui suivent, on se propose de démontrer le théorème de Cauchy : si $G$ est un groupe fini de cardinal $n$ et $p$ est un nombre premier divisant $n$, alors $G$ possède un élément d'ordre $p$.
\begin{ex}{Un lemme}
	On suppose $G$ abélien. Soient $x$ et $y$ deux éléments de $G$ d'ordres respectifs $a$ et $b$ premiers entre eux. Montrer que $xy$ est d'ordre $ab$.
\end{ex}
\begin{ex}{Preuve du théorème dans le cas abélien}
	Soit $G$ un groupe fini abélien d'ordre $n$. Soit $p$ un nombre premier divisant $n$. On note $x_1,\ldots,x_n$ les éléments de $G$. Pour $i$ variant de $1$ à $n$, on note $\omega_i$ l'ordre de $x_i$. On définit l'application : $$\begin{array}{lrcl}
		f:&\Z/\omega_1\Z\times\cdots\times\Z/\omega_n\Z&\longrightarrow&G \\
		&(\overline{a_1},\ldots,\overline{a_n})&\longmapsto&x_1^{a_1}\cdots x_n^{a_n}
	\end{array}$$ Justifier que $f$ est bien définie et est un morphisme de groupes. En déduire que $n$ divise $\displaystyle\prod_{i=1}^{n}\omega_i$ puis que $p$ divise l'un des $\omega_i$. Conclure.
\end{ex}
\begin{ex}{Une première application}
	On suppose que $G$ est un groupe abélien d'ordre $p_1\cdots p_r$ où les $p_i$ sont des nombres premiers deux à deux distincts. Montrer qu $G$ est cyclique.
\end{ex}
\begin{ex}{Preuve du théorème dans le cas général}
	On ne suppose plus nécessairement $G$ abélien. On pose $E=\{(x_1,\ldots,x_p)\in G_p\ \mid\ x_1\cdots x_p=e\}$. On définit dans $E$ la relation binaire $\mathcal{R}$ de la façon suivante : $$(x_1,\ldots,x_p)\mathcal{R}(y_1,\ldots,y_p)\iff \exists k\in\llbracket0,p-1\rrbracket,\ \begin{cases}
		\forall i\in\llbracket1,p-k\rrbracket,\ y_i=x_{k+i}\\ \forall i\in\llbracket p-k+1,p\rrbracket,\ y_i=x_{i-p+k}
	\end{cases}$$
	\begin{enumerate}
		\item Justifier que $\card(E)=n^{p-1}$.
		\item Vérifier que $\mathcal{R}$ est une relation d'équivalence. Quelle est la classe d'équivalence de l'élément $(e,\ldots,e)$ ?
		\item Soit $X\in E$. Prouver l'équivalence : $\card(\dot{X})=1\iff\exists x\in G,\ X=(x,\ldots,x)\ \text{avec}\ x^p=e$. Montrer que si $\card(\dot{X})\ne 1$, alors $\card(\dot{X})=p$.
		\item Soit $H=\{x\in G\ \mid x^p=e\}$. Déduire de ce qui précède que $\card(H)$ est un multiple de $p$.
		\item Conclure.
	\end{enumerate}
\end{ex}
\begin{proof}
	Pour la 4, passer aux cardinaux le fait que les classes d'équivalence forment une partition.
\end{proof}
\begin{ex}{Une deuxième application}
	\begin{enumerate}
		\item Démontrer que si $p$ est premier, il existe deux structures de groupe d'ordre d'ordre $p^2$ : la structure cyclique et celle de $(\Z/p\Z)^2$.
		\item Démontrer que si $p$ est premier, il existe deux structures de groupe d'ordre $2p$ : la structure cyclique et celle du groupe diédral $D_p$.
		\item Étudier les structures de groupe d'ordre 2024  et 2026.
	\end{enumerate}
\end{ex}
\begin{ex}{Un théorème de Frobenius}
	Soient $G$ un groupe fini et $H$ un sous-groupe fini de $G$. On note $p$ le plus petit diviseur premier du cardinal de $G$. On suppose que $H$ est d'indice fini $p$ dans $G$, c'est-à-dire $\card(G)=p\times\card(H)$.
	\begin{enumerate}
		\item On note $(G/H)_g$ l'ensemble des classes à gauche selon $H$. Démontrer que l'application $(x,aH)\mapsto xaH$ est une action du groupe $G$ sur $(G/H)_g$.
		\item Démontrer que $\{\sigma_x\ \mid \ x\in G\}$ est de cardinal $p$.
		\item En déduire que $H$ est distingué dans $G$.
	\end{enumerate}
\end{ex}
\subsection{Théorème de l'élément primitif}
\begin{ex}{Autour des ordres et une formule d'Euler}
	\begin{enumerate}
		\item Soit $a\in\llbracket0,n-1\rrbracket$. Montrer que l'ordre de $\overline{a}$ dans le groupe additif $\Z/n\Z$ est égal à $\displaystyle\frac{n}{a\wedge n}$.
		\item En regroupant les éléments selon leur ordre, démontrer que si $n\in\N\st$, alors on a : $$n=\sum_{d\mid n}\varphi(d)$$
	\end{enumerate}
\end{ex}
\begin{ex}{Théorème de l'élément primitif}
	Soit $\K$ un corps et $G$ un sous-groupe fini de $\K\st$ de cardinal $n$.
	\begin{enumerate}
		\item Si $d$ est un diviseur de $n$, on note $\psi(d)$ le nombre d'éléments de $G$ d'ordre $d$. Justifier que $n=\displaystyle\sum_{d\mid n}\psi(d)$.
		\item Soit $d$ un diviseur de $n$. On suppose qu'il existe un élément de $G$ d'ordre $d$. Justifier que le polynôme $X^d-1$ possède exactement $d$ racines dans $\K$. En déduire que $G$ admet exactement $\varphi(d)$ éléments d'ordre $d$.
		\item En utilisant les questions précédentes et l'exercice précédent, montrer que pour tout $d$ diviseur de $n$, on a : $\psi(d)=\varphi(d)$. En déduire que $G$ est cyclique.
	\end{enumerate}
\end{ex}
\begin{rmq}
	Ce résultat constitue le théorème de l'élément primitif. En l'appliquant au corps $\K=\Z/p\Z$ pour $p$ premier et au groupe $G=\K\st$, on en déduit que le groupe multiplicatif $(\Z/p\Z)\st$ est cyclique. Un générateur de ce groupe s'appelle un élément primitif modulo $p$.
\end{rmq}
\begin{ex}{Détermination pratique d'éléments primitifs modulo $p$}
	Déterminer un élément primitif modulo $p$ pour $p=5$ puis pour $p=7$.
\end{ex}
\subsection{Structure du groupe multiplicatif des inversibles de $\Z/n\Z$}
\begin{ex}{Un lemme utile}
	Soit $(G,\cdot)$ un groupe abélien. Soient $x$ et $y$ deux éléments de $G$ d'ordres respectifs $m$ et $n$ premiers entre eux. Montrer que l'ordre de $x\cdot y$ est $mn$.
\end{ex}
\begin{ex}{Structure des $(\Z/p^n\Z)^\times$ pour $p$ premier}
	Soit $p$ un nombre premier différent de $2$. On se place dans le groupe multiplicatif $G=(\Z/p^n\Z)^\times$.
	\begin{enumerate}
		\item Montrer par récurrence que si $k$ est un entier naturel supérieur ou égal à $2$, alors $$(1+p)^{p^{k-2}}\ \equiv\ 1+p^{k-1}\ [p^k]$$ 
		\item En déduire que $$(1+p)^{p^{k-1}}\ \equiv\ 1\ [p^k]$$
		\item En déduire que $\overline{1+p}$ est d'ordre $p^{n-1}$ dans $G$.
		\item Quel est le cardinal de $G$ ?
		\item En utilisant un générateur du groupe multiplicatif $(\Z/p\Z)\st$ (dont l'existence est assurée par la partie précédente), montrer que $G$ est cyclique.
	\end{enumerate}
\end{ex}
\begin{ex}{Étude du cas particulier $(\Z/2^n\Z)^\times$}
	On suppose ici $p=2$ et $n\geqslant 2$. On étudie donc le groupe multiplicatif $G=(\Z/2^n\Z)^\times$.
	\begin{enumerate}
		\item Montrer que si $k$ est un entier naturel, alors $$5^{2^k}\ \equiv\ 1+2^{k+2}\ [2^{k+3}]$$ Indication : on pourra procéder par récurrence.
		\item En déduire que $\overline{5}$ est d'ordre $2^{n-2}$ dans $G$.
		\item Montrer que $\overline{-1}$ n'appartient pas au groupe engendré par $\overline{5}$ dans $G$.
		\item En déduire que, pour $n$ supérieur ou égal à $3$, $G$ est isomorphe au groupe additif $(\Z/2^{n-2}\Z)\times(\Z/2\Z)$.
		\item En déduire que $G$ n'est pas cyclique.
	\end{enumerate}
\end{ex}
\begin{ex}{Dévissage de $(\Z/n\Z)^\times$}
	\begin{enumerate}
		\item Soient $G$ et $G'$ deux groupes multiplicatifs cycliques d'ordre respectifs $m$ et $n$. On munit $G\times G'$ de la structure de groupe produit. Montrer que $G\times G'$ est cyclique si, et seulement si, $m$ et $n$ sont premiers entre eux.
		\item Si $m$ et $n$ sont deux entiers premiers entre eux, rappeler rapidement pourquoi les groupes multiplicatifs $(\Z/mn\Z)^\times$ et $(\Z/m\Z)^\times\times(\Z/n\Z)^\times$ sont isomorphes.
		\item Soit $n\geqslant 2$ un entier que l'on décompose en facteurs premiers sous la forme $\displaystyle\prod_{i=1}^{r}p_i^{a_i}$. Montrer que $(\Z/n\Z)^\times$ est isomorphe à $$\prod_{i=1}^{r}(\Z/p_i^{a_i}\Z)^\times$$
		\item Démontrer que le groupe multiplicatif $(\Z/n\Z)^\times$ est cyclique si, et seulement si, $n=2$, $n=4$, $n=p^a$ ou $n=2p^a$, où $p$ est un nombre premier impair.
	\end{enumerate}
\end{ex}

\subsection{Construction de corps finis}
Pour $p$ premier, on notera $\F_p$ le corps $\Z/p\Z$.
\begin{ex}{Quotient de $\K[X]$ par un idéal principal}
	Soit $\K$ un corps et $P$ un polynôme non constant de $\K[X]$. On introduit sur $\K[X]$ la relation binaire $\mathcal{R}$ définie par : $A\mathcal{R}B\iff \exists Q\in\K[X],\ P=Q(A-B)$.
	\begin{enumerate}
		\item Vérifier que $\mathcal{R}$ est une relation d'équivalence sur $\K[X]$ compatible avec les lois $+$ et $\cdot$. \\
		On note $\overline{A}$ la classe d'équivalence du polynôme $A$, c'est-à-dire l'ensemble $A+P\K[X]$. On note $\K[X]/(P)$ l'ensemble quotient, dans lequel on définit naturellement les lois quotient. Ainsi, $\K[X]/(P)$ est muni d'une structure d'anneaux et la projection canonique est un morphisme d'anneaux.
		\item Démontrer que l'application $j:\K\to\K[X]/(P)$ définie par $j(u)=\overline{u}$ est injective. Ce morphisme d'anneaux injectif permet d'identifier $\K$ et on image $j(\K)$.
		\item Démontrer que $\K[X]/(P)$ est un corps si, et seulement si, $P$ est irréductible dans $\K[X]$. Indication : on pourra utiliser la réciproque de l'identité de Bézout. \\
		On supposera désormais que $P$ est irréductible et on notera $\K'$ ce corps. $j$ permet alors d'identifier $\K$ à un sous-corps de $\K'$. Ainsi, on peut abusivement considérer que $\K[X]\subset\K'[X]$.
		\item On pose $\alpha=\overline{X}$. Justifier que $P(\alpha)=0$. On en déduit que le polynôme $P$, qui était irréductible dans $\K[X]$, ne l'est plus dans $\K'[X]$ puisque $\alpha$ en est une racine.
		\item En déduire que l'on peut construire une extension du corps $\K$, c'est-à-dire un corps $\L$ contenant un sous-corps isomorphe à $\K$, dans laquelle tout polynôme est scindé.
	\end{enumerate}
\end{ex}
\begin{rmq}
	Une telle extension de corps est appelée corps de décomposition de $P$.
\end{rmq}
\begin{ex}{Nombre de polynômes irréductibles à coefficients dans $\F_p$}
	Soit $p$ un nombre premier. On s'intéresse au polynômes irréductibles de l'anneau $\F_p[X]$. on note $m_r(p)$ le nombre de polynômes irréductibles unitaires de degré $r$ dans $\F_p[X]$.
	\begin{enumerate}
		\item Démontrer que $m_1(p)=p$ et $m_2(p)=\displaystyle\frac{p(p-1)}{2}$. Indication : pour la deuxième égalité, on pourra dénombrer les polynômes réductibles de degré $2$.
		\item Démontrer que dans $\F_p[X]$, le polynôme $X^{p^n}-X$ est exactement le produit de tous les polynômes irréductibles unitaires dont le degré divise $n$. Vérifier ce résultat lorsque $p=2$ et $n=3$.
		\item Déduire de la question précédente que $p^n=\displaystyle\sum_{r\mid n}rm_r(p)$.
		\item En déduire l'encadrement suivant : $$\frac{p^n-p{E(n/2)+1}}{n}\leqslant m_n(p)\leqslant\frac{p^n}{n}$$
		\item En déduire qu'il existe dans $\F_p[X]$ des polynômes irréductibles de tout degré.
	\end{enumerate}
\end{ex}
\begin{ex}{Construction d'un corps fini de cardinal $p^n$}
	Soit $p$ un nombre premier. Soit $n\in\N\st$. Soit $P$ un polynôme irréductible de degré $n$ à coefficients dans $\F_p$. On construit comme dans la section précédente un corps $\K'=\F_p[X]/(P)$.
	\begin{enumerate}
		\item Si $k\in\N$, on note $E_k$ l'ensemble des polynômes à coefficients dans $\F_p$ de degré inférieur ou égal à $k$. Montrer que l'application $s:E_{n-1}\to\K'$ définie par $s(A)=\overline{A}$ est un isomorphisme d'espaces vectoriels.
		\item En déduire que le cardinal de $\K'$ est égal à $p^n$.
		\item On choisit $p=2$ et $P=X^2+X+1$. Vérifier que $P$ est irréductible dans $\F_2[X]$. On note $\F_4$ le corps $\F_2[X]/(P)$. En notant $\alpha=\overline{A}$, montrer que $\F_4=\{0,1,\alpha,\alpha+1\}$. Justifier que $\alpha^2=\alpha+1$. Dresser la table de multiplication de ce corps.
		\item On choisit $p=3$ et $P=X^2+1$. Justifier que $P$ est irréductible dans $\F_3[X]$. Donner la liste des $9$ éléments du corps $\F_3[X]/(P)$ que l'on note $\F_9$. Indiquer comment construire la table de multiplication de ce corps.
		\item Montrer que le cardinal d'un corps fini est nécessairement une puissance d'un nombre premier.
	\end{enumerate}
\end{ex}
\begin{rmq}
	On a montré que pour $p$ premier et $n\in\N\st$, il existe un corps de cardinal $p^n$. On peut prouver que ce corps de cardinal $p^n$ est unique à isomorphisme près.
\end{rmq}


\end{document}